\documentclass[aps,prb,twocolumn,groupedaddress,amsmath,amssymb,
               superscriptaddress,nofootinbib,floatfix]{revtex4-2}

\usepackage{bm}
\usepackage{graphicx}
\usepackage{booktabs}
\usepackage{xcolor}
\usepackage{hyperref}
\usepackage{amsthm}
\usepackage{array}

\hypersetup{
  colorlinks=true,
  citecolor=blue!55!black,
  linkcolor=blue!55!black,
  urlcolor=blue!55!black,
  pdftitle={A Finite-Volume Infrared Theorem for Classical O(n) Spin
            Systems: Absence of Continuous-Symmetry Long-Range Order
            in Two Dimensions}
}

\theoremstyle{plain}
\newtheorem{theorem}{Theorem}
\newtheorem{lemma}{Lemma}
\newtheorem{corollary}{Corollary}
\theoremstyle{definition}
\newtheorem{assumption}{Assumption}

\newcommand{\BZ}{\mathrm{BZ}}
\newcommand{\dd}{\mathrm{d}}
\newcommand{\cL}{c_L}
\newcommand{\cinf}{c_\infty}
\newcommand{\rhoLog}{\rho_{\log}}
\newcommand{\kmin}{k_{\min}}
\newcommand{\Torus}{\mathbb T}
\newcommand{\Exp}{\mathbb E}

\begin{document}

\title{A Finite-Volume Infrared Theorem for Classical
       \texorpdfstring{$O(n)$}{O(n)} Spin
       Systems: Absence of Continuous-Symmetry Long-Range Order in
       Two Dimensions}

\author{Authors withheld for internal draft}
\affiliation{Affiliations to be finalized}
\date{\today}

\begin{abstract}
We prove a finite-volume infrared theorem for classical hard-spin
$O(n)$ models, $n\ge2$, with translation-invariant bilinear ferromagnetic
interactions.  For a torus of volume $V=L^d$, let
$E_L(q)=\sum_RJ_L(R)[1-\cos(q\cdot R)]$.  If $E_L$ converges uniformly to
a continuous nonnegative kernel $E_\infty$ and
\[
 I(h)=\int_{\BZ}\frac{\dd^dq}{(2\pi)^d}
 \frac{1}{h+E_\infty(q)}
\]
diverges as $h\downarrow0$, then
\[
 \lim_{h\downarrow0}\limsup_{L\to\infty}m_{L,h}=0,
 \qquad
 \lim_{L\to\infty}\langle|\bm M_L|^2\rangle_{L,0}=0
\]
at every finite temperature.  The proof uses divergence-free rotation
generators on $S^{n-1}$.  The exact pair-energy second variation is a
two-component internal-plane projection, and summing the $n-1$ transverse
channels yields
\[
 1\ge(n-1)T m_{L,h}^2
 \frac1V\sum_q\frac1{h+E_L(q)}.
\]
An exponential-tilt lemma then connects this field-selected estimate to the
zero-field squared magnetization measured in simulations.  In two
dimensions the theorem covers finite-second-moment interactions and the
marginal $J(R)\propto |R|^{-4}$ model, including the minimum-image,
size-normalized torus sequence.  At the marginal point we further derive the
rigorous asymptotic finite-size envelope
\[
 \langle|\bm M_L|^2\rangle_{L,0}
 =O_{n,T}\!\left((\log\log L)^{-1/2}\right),
\]
which applies directly to the zero-field Monte Carlo observable but is not
claimed to be its exact decay rate.  A deterministic lattice audit through
$L=8192$ reproduces the exact logarithmic kernel coefficient to better than
$9.1\times10^{-7}$ along four momentum paths.  Public-data diagnostics for
the $n=2$ application explain why a zero asymptote can look ordered over
accessible sizes.  The theorem does not apply to arbitrary continuous
fields or to long-range kernels with a finite infrared integral, and it does
not determine the correlation law of the remaining no-order phase.
\end{abstract}

\maketitle

\section{Introduction}

The absence of continuous-symmetry breaking in low-dimensional systems is
usually associated with the Mermin--Wagner theorem
\cite{MerminWagner1966}.  Its classical formulation is not merely a
large-spin limit of a quantum commutator argument: a purely classical
Bogoliubov inequality follows from integration by parts in the Gibbs
measure \cite{Mermin1967}.  For long-range interactions, the decisive object
is the excitation kernel
\begin{equation}
 E(q)=\sum_RJ(R)[1-\cos(q\cdot R)]                      \label{eq:intro-kernel}
\end{equation}
and the infrared integral of its reciprocal
\cite{Bruno2001,Messager1984,Ioffe2002}.

Two qualifications are often compressed into the phrase ``continuous
theory.''  First, the existence of a continuous global symmetry does not
specify the single-site space, the interaction, or the small-momentum
kernel.  Second, a theorem that rules out a field-selected one-point
function does not by itself display the finite-volume route to
\begin{equation}
 \langle|\bm M_L|^2\rangle_{L,0},                       \label{eq:mc-M2}
\end{equation}
the zero-field observable extrapolated in Monte Carlo work.  At a marginal
infrared boundary, this distinction matters because the thermodynamic
limit can be approached more slowly than any visible power of $L$.

The motivating example is the two-dimensional interaction
\begin{equation}
 J(r)\propto r^{-(2+\sigma)}.                           \label{eq:power-law}
\end{equation}
Recent simulations of the XY case reported ferromagnetic order up to
$\sigma=2$ \cite{Xiao2025,Yao2025}.  At that endpoint,
$J(r)\propto r^{-4}$ and
\begin{equation}
 E(q)\sim\rhoLog q^2\log(1/q).                          \label{eq:intro-marginal}
\end{equation}
The reciprocal integral diverges only as $\log\log L$.  The resulting
finite-size sequence can remain close to a plateau at sizes far beyond
current simulations.

This article separates three statements.  The first is an abstract
finite-volume theorem for hard-spin $O(n)$ systems in arbitrary spatial
dimension.  The second verifies its hypotheses for two-dimensional
short-range and marginal $1/r^4$ interactions.  The third uses deterministic
kernel calculations and public simulation summaries to examine finite-size
identifiability.  The numerical work is not a premise of the theorem.

The extension from XY to $O(n)$ is not obtained by replacing the number
$2$ with $n$.  Angular integration by parts must be replaced by rotation
calculus on $S^{n-1}$; the pair second variation is a selected-plane
projection rather than the full scalar product; and all $n-1$ transverse
spectral budgets must be combined without overcounting.  We give those steps
explicitly.  They close for unit-vector spins with bilinear ferromagnetic
couplings.  They do not establish a theorem for unbounded fields, arbitrary
target manifolds, frustrated interactions, or quantum models.

\section{Model, hypotheses, and result}
\label{sec:model}

Let
\begin{equation}
 \Torus_L^d=(\mathbb Z/L\mathbb Z)^d,\qquad V=L^d,
\end{equation}
and assign a unit vector
\begin{equation}
 \bm S_x=(S_x^1,\ldots,S_x^n)\in S^{n-1},
 \qquad n\ge2,                                         \label{eq:hard-spin}
\end{equation}
to every site.  The single-site measure is the normalized
$SO(n)$-invariant surface measure.  We consider
\begin{equation}
 H_{L,h}
 =
 -\frac12\sum_{x\in\Torus_L^d}\sum_R
 J_L(R)\,\bm S_x\cdot\bm S_{x+R}
 -h\sum_xS_x^1,                                        \label{eq:H}
\end{equation}
with
\begin{equation}
 J_L(R)=J_L(-R)\ge0,\qquad J_L(0)=0.                   \label{eq:J-assumption}
\end{equation}
The factor $1/2$ compensates the directed-pair sum, including
self-inverse torus displacements.

At $T=\beta^{-1}>0$, define
\begin{equation}
 \bm M_L=\frac1V\sum_x\bm S_x,\qquad
 m_{L,h}=\langle M_L^1\rangle_{L,h},                   \label{eq:magnetizations}
\end{equation}
and, for torus momentum $q$,
\begin{equation}
 E_L(q)=\sum_RJ_L(R)[1-\cos(q\cdot R)].                \label{eq:E_L}
\end{equation}
The fixed integer representatives of $R$ give a continuous periodic
interpolation of $E_L$ on the Brillouin zone.

\begin{assumption}[Uniform kernel limit]
\label{ass:uniform}
There is a continuous nonnegative periodic function $E_\infty$ such that
\begin{equation}
 \|E_L-E_\infty\|_\infty\longrightarrow0.              \label{eq:uniform}
\end{equation}
\end{assumption}

\begin{assumption}[Regulated infrared divergence]
\label{ass:ir}
The integral
\begin{equation}
 I(h)=\int_{\BZ}\frac{\dd^dq}{(2\pi)^d}
 \frac1{h+E_\infty(q)}                                 \label{eq:Ih}
\end{equation}
satisfies $I(h)\to\infty$ as $h\downarrow0$.
\end{assumption}

\begin{theorem}[Finite-volume hard-spin $O(n)$ criterion]
\label{thm:main}
Let $n\ge2$, $0<T<\infty$, and suppose
Eqs.~\eqref{eq:J-assumption}, \eqref{eq:uniform}, and
\eqref{eq:Ih} hold.  Then
\begin{align}
 \lim_{h\downarrow0}\limsup_{L\to\infty}m_{L,h}&=0,
                                                               \label{eq:field-result}\\
 \lim_{L\to\infty}
 \langle|\bm M_L|^2\rangle_{L,0}&=0.                  \label{eq:zero-result}
\end{align}
At finite volume,
\begin{equation}
 1\ge
 (n-1)T\,m_{L,h}^{\,2}
 \frac1V\sum_{q\in\mathcal Q_L}
 \frac1{h+E_L(q)}.                                     \label{eq:master-bound}
\end{equation}
\end{theorem}

Table~\ref{tab:scope} records the logical boundary of the result.  The
abstract theorem is not restricted to $d=2$; its two-dimensional content
comes from interactions for which Assumption~\ref{ass:ir} can be verified.

\begin{table*}[t]
\caption{\label{tab:scope}
Hypotheses used by Theorem~\ref{thm:main}.  The right column lists changes
that require a new argument rather than a change of notation.}
\centering
\renewcommand{\arraystretch}{1.18}
\resizebox{0.98\textwidth}{!}{%
\begin{tabular}{lll}
\toprule
Object & Assumption in the theorem & Outside the proved scope\\
\midrule
Single-site space &
\shortstack[l]{$S^{n-1}$, finite $n\ge2$,\\normalized invariant measure} &
\shortstack[l]{$n=1$, unbounded soft spins,\\a generic target manifold}\\
\addlinespace[3pt]
Interaction &
\shortstack[l]{translation-invariant bilinear pair energy,\\$J_L(R)\ge0$} &
\shortstack[l]{frustration, sign changes,\\random or multibody terms}\\
\addlinespace[3pt]
Thermodynamic input &
\shortstack[l]{uniform kernel limit and divergent\\regulated reciprocal integral} &
\shortstack[l]{size sequences without a kernel limit,\\or finite infrared integral}\\
\addlinespace[3pt]
Conclusion &
\shortstack[l]{no vector LRO in a field-selected state;\\no zero-field $M_L^2$ intercept} &
\shortstack[l]{no assertion about the exact\\correlation law or topological sector}\\
\bottomrule
\end{tabular}
}
\end{table*}

\section{Rotation calculus on the sphere}
\label{sec:rotation}

For each transverse component $a=2,\ldots,n$, let
\begin{equation}
 \mathcal L_x^{1a}
 =
 S_x^1\frac{\partial}{\partial S_x^a}
 -
 S_x^a\frac{\partial}{\partial S_x^1}.                 \label{eq:generator}
\end{equation}
Equation~\eqref{eq:generator} denotes the derivative along the rotation
orbit in the $(1,a)$ plane.  If $R_{1a}(t)\in SO(n)$ is that rotation, then
\begin{equation}
 \mathcal L^{1a}f(\bm S)
 =
 \left.\frac{\dd}{\dd t}f(R_{1a}(t)\bm S)\right|_{t=0}. \label{eq:orbit-derivative}
\end{equation}
The two relevant component identities are
\begin{equation}
 \mathcal L^{1a}S^1=-S^a,\qquad
 \mathcal L^{1a}S^a=S^1.                               \label{eq:component-rotation}
\end{equation}

\begin{lemma}[Invariant-measure integration by parts]
\label{lem:sphere-ibp}
For a smooth function $f$ on $S^{n-1}$,
\begin{equation}
 \int_{S^{n-1}}\mathcal L^{1a}f\,\dd\omega=0.           \label{eq:sphere-ibp}
\end{equation}
\end{lemma}

\begin{proof}
Rotation invariance of $\omega$ gives
\begin{equation}
 \int f(R_{1a}(t)\bm S)\,\dd\omega(\bm S)
 =
 \int f(\bm S)\,\dd\omega(\bm S).
\end{equation}
Differentiate at $t=0$.  This proof also identifies the geometric
subtlety: $\mathcal L^{1a}$ is a divergence-free Killing field.  An ambient
coordinate derivative not tangent to the sphere would produce additional
measure terms.
\end{proof}

Define
\begin{equation}
 D_q^a=\sum_xe^{iq\cdot x}\mathcal L_x^{1a},\qquad
 A_q^a=\sum_xe^{-iq\cdot x}S_x^a.                      \label{eq:D-A}
\end{equation}
Applying Lemma~\ref{lem:sphere-ibp} to the product Gibbs measure gives
\begin{align}
 \langle D_q^aF\rangle
 &=\beta\langle F D_q^aH\rangle,                       \label{eq:gibbs-ibp-one}\\
 \langle D_{-q}^aD_q^aH\rangle
 &=\beta\langle|D_q^aH|^2\rangle\ge0.                 \label{eq:gibbs-ibp-two}
\end{align}
The nonnegativity in Eq.~\eqref{eq:gibbs-ibp-two} is a Gibbs-averaged
identity.  The configuration-level second variation need not be positive.

\section{Finite-volume infrared estimate}
\label{sec:finite-proof}

\subsection{Classical Bogoliubov inequality}

Set $F=A_q^a$ in Eq.~\eqref{eq:gibbs-ibp-one}.  Cauchy--Schwarz and
Eq.~\eqref{eq:gibbs-ibp-two} yield
\begin{equation}
 \langle|A_q^a|^2\rangle
 \langle D_{-q}^aD_q^aH\rangle
 \ge
 T|\langle D_q^aA_q^a\rangle|^2.                      \label{eq:classical-bog}
\end{equation}
The numerator is independent of $q$ and $a$:
\begin{equation}
 D_q^aA_q^a=\sum_xS_x^1,\qquad
 \langle D_q^aA_q^a\rangle=Vm_{L,h}.                  \label{eq:numerator}
\end{equation}

\subsection{Exact pair-energy second variation}

Let $s,t\in S^{n-1}$ denote a single pair and fix a transverse component
$a$.  Define
\begin{equation}
 B=s_1t_a-s_at_1,\qquad C=s_1t_1+s_at_a.               \label{eq:B-C}
\end{equation}
For the generators acting on $s$ and $t$,
\begin{equation}
\begin{array}{ll}
 \mathcal L_s(s\cdot t)=B,&
 \mathcal L_t(s\cdot t)=-B,\\[1mm]
 \mathcal L_sB=-C,&
 \mathcal L_tB=C.
\end{array}                                             \label{eq:pair-algebra}
\end{equation}
Consequently, for complex $u,v$,
\begin{multline}
 (\bar u\mathcal L_s+\bar v\mathcal L_t)
 (u\mathcal L_s+v\mathcal L_t)(-J\,s\cdot t)\\
 =
 J|u-v|^2(s_1t_1+s_at_a).                              \label{eq:pair-variation}
\end{multline}
For $u=e^{iq\cdot x}$ and $v=e^{iq\cdot(x+R)}$,
$|u-v|^2=2[1-\cos(q\cdot R)]$.  Combining this with the directed-pair
factor in Eq.~\eqref{eq:H} gives
\begin{align}
 D_{-q}^aD_q^aH_{L,h}
 ={}&h\sum_xS_x^1 \nonumber\\
 &+\sum_{x,R}J_L(R)[1-\cos(q\cdot R)]\nonumber\\
 &\quad\times
 (S_x^1S_{x+R}^1+S_x^aS_{x+R}^a).                    \label{eq:exact-denominator}
\end{align}

The pair factor in Eq.~\eqref{eq:exact-denominator} is not
$\bm S_x\cdot\bm S_{x+R}$.  It is the projection onto the internal
$(1,a)$ plane.  Cauchy--Schwarz in that plane gives
\begin{equation}
 S_x^1S_y^1+S_x^aS_y^a\le1.                            \label{eq:projection-bound}
\end{equation}
It may also be negative.  Ferromagnetic positivity and
$S_x^1\le1$ therefore imply the upper bound
\begin{equation}
 \langle D_{-q}^aD_q^aH_{L,h}\rangle
 \le V[h+E_L(q)].                                      \label{eq:denominator-upper}
\end{equation}
Equations~\eqref{eq:classical-bog}, \eqref{eq:numerator}, and
\eqref{eq:denominator-upper} give, for every transverse channel,
\begin{equation}
 \langle|A_q^a|^2\rangle
 \ge
 \frac{TVm_{L,h}^2}{h+E_L(q)}.                         \label{eq:channel-bound}
\end{equation}

\subsection{Summing the transverse budgets}

Parseval's identity gives
\begin{align}
 \sum_{a=2}^n\sum_q|A_q^a|^2
 &=
 V\sum_x\sum_{a=2}^n(S_x^a)^2 \nonumber\\
 &=
 V\sum_x[1-(S_x^1)^2]\le V^2.                         \label{eq:parseval-budget}
\end{align}
Summing Eq.~\eqref{eq:channel-bound} over $q$ and over all $n-1$
transverse components proves Eq.~\eqref{eq:master-bound}.

This is the origin of the factor $n-1$.  It is not obtained by multiplying
one inequality while keeping the same left-hand side; the combined
transverse spectral weight has the single shared budget
\eqref{eq:parseval-budget}.  For $n=2$ there is one transverse component and
the formula reduces exactly to the XY bound.

\section{Thermodynamic limit and the zero-field bridge}
\label{sec:bridge}

Fix $h>0$.  Assumption~\ref{ass:uniform} and
\begin{equation}
 \left|\frac1{h+x}-\frac1{h+y}\right|
 \le\frac{|x-y|}{h^2},\qquad x,y\ge0,                  \label{eq:reciprocal-lipschitz}
\end{equation}
imply uniform convergence of the reciprocal kernels.  The momentum average
is therefore a Riemann sum:
\begin{equation}
 \frac1V\sum_q\frac1{h+E_L(q)}
 \longrightarrow I(h).                                \label{eq:riemann}
\end{equation}
Equation~\eqref{eq:master-bound} yields
\begin{equation}
 \limsup_{L\to\infty}m_{L,h}^2
 \le\frac1{(n-1)T I(h)}.                               \label{eq:regulated-bound}
\end{equation}
Assumption~\ref{ass:ir} then proves Eq.~\eqref{eq:field-result}.

It remains to reach Eq.~\eqref{eq:mc-M2}.  The following scalar statement is
independent of the spin model.

\begin{lemma}[Bounded symmetric exponential tilt]
\label{lem:tilt}
Let $X\in[-1,1]$ have a distribution symmetric under $X\mapsto-X$.  For
$t\ge0$, let
\begin{equation}
 \Exp_t[f(X)]
 =
 \frac{\Exp[f(X)e^{tX}]}{\Exp[e^{tX}]}.
\end{equation}
Then
\begin{equation}
 \Exp_t[X]\ge\tanh(t)\,\Exp[X^2].                      \label{eq:tilt}
\end{equation}
\end{lemma}

\begin{proof}
Symmetry gives
\begin{equation}
 \Exp_t[X]
 =
 \frac{\Exp[X\sinh(tX)]}{\Exp[\cosh(tX)]}.
\end{equation}
For $0\le x\le1$, concavity of $\tanh$ gives
$\tanh(tx)\ge x\tanh(t)$, and hence
\begin{equation}
 x\sinh(tx)\ge x^2\tanh(t)\cosh(tx).
\end{equation}
Both $x^2$ and $\cosh(tx)$ increase with $x$ on $[0,1]$; their covariance
under the law of $|X|$ is nonnegative.  Thus
\begin{equation}
 \Exp[X^2\cosh(tX)]
 \ge\Exp[X^2]\Exp[\cosh(tX)],
\end{equation}
which proves Eq.~\eqref{eq:tilt}.
\end{proof}

At $h=0$, $SO(n)$ invariance makes the magnetization second-moment tensor
isotropic:
\begin{equation}
 \langle M_L^\alpha M_L^\gamma\rangle_{L,0}
 =
 \frac{\delta_{\alpha\gamma}}{n}
 \langle|\bm M_L|^2\rangle_{L,0}.                      \label{eq:isotropic-M2}
\end{equation}
A rotation by $\pi$ in the $(1,2)$ plane makes the distribution of
$X=M_L^1$ symmetric.  The field changes the zero-field Gibbs weight by
$e^{\beta hVX}$.  Lemma~\ref{lem:tilt}, with $t=\beta hV$, therefore gives
\begin{equation}
 m_{L,h}
 \ge
 \frac1n\tanh(\beta hV)
 \langle|\bm M_L|^2\rangle_{L,0}.                      \label{eq:bridge}
\end{equation}
At fixed $h>0$, the hyperbolic tangent tends to one, so
\begin{equation}
 \limsup_{L\to\infty}
 \langle|\bm M_L|^2\rangle_{L,0}
 \le n\limsup_{L\to\infty}m_{L,h}.                    \label{eq:bridge-limit}
\end{equation}
Taking $h\downarrow0$ only after the thermodynamic limit and using
Eq.~\eqref{eq:regulated-bound} proves Eq.~\eqref{eq:zero-result}.

\section{Two-dimensional consequences}
\label{sec:2d}

\subsection{Finite second moment}

For a fixed interaction satisfying
\begin{equation}
 \sum_RJ(R)|R|^2<\infty,                               \label{eq:second-moment}
\end{equation}
the inequality $1-\cos u\le u^2/2$ gives
\begin{equation}
 E_\infty(q)
 \le\frac{|q|^2}{2}\sum_RJ(R)|R|^2
 =C|q|^2.                                              \label{eq:q2-bound}
\end{equation}
In $d=2$, the integral of $1/(h+Cq^2)$ diverges logarithmically as the field
is removed.  Theorem~\ref{thm:main} therefore recovers the standard
short-range no-order result for every finite $n\ge2$.

\subsection{Marginal \texorpdfstring{$1/r^4$}{1/r4} interaction}

Let
\begin{equation}
 J_\infty(R)=\frac{A}{|R|^4},
 \qquad R\in\mathbb Z^2\setminus\{0\}.                 \label{eq:r4}
\end{equation}
Split the kernel at $|R|=|q|^{-1}$.  In the near region,
\begin{align}
 E_{<}(q)
 &\le
 \frac{A|q|^2}{2}
 \sum_{0<|R|\le|q|^{-1}}\frac1{|R|^2} \nonumber\\
 &\le C_1|q|^2\log\frac{C_2}{|q|}.                    \label{eq:near-bound}
\end{align}
In the far region,
\begin{equation}
 E_{>}(q)
 \le2A\sum_{|R|>|q|^{-1}}\frac1{|R|^4}
 \le C_3|q|^2.                                         \label{eq:far-bound}
\end{equation}
Hence
\begin{equation}
 E_\infty(q)
 \le C|q|^2\log\frac{C'}{|q|},                         \label{eq:marginal-upper}
\end{equation}
and
\begin{equation}
 \int_0^\delta
 \frac{k\,\dd k}{k^2\log(C'/k)}
 =\infty.                                              \label{eq:loglog-divergence}
\end{equation}
The divergence is only $\log\log$, but it is sufficient for
Theorem~\ref{thm:main}.

The leading square-lattice coefficient follows from
\begin{equation}
 \sum_{0<|R|\le M}
 \frac{R_iR_j}{|R|^4}
 =
 \pi\delta_{ij}\log M+O(1),                            \label{eq:tensor-sum}
\end{equation}
which gives
\begin{equation}
 E_\infty(q)
 =
 \frac{\pi A}{2}|q|^2\log\frac1{|q|}
 +O(|q|^2).                                            \label{eq:exact-kernel}
\end{equation}
The exact coefficient is useful for convention checks; the upper bound
\eqref{eq:marginal-upper} is enough for the theorem.

\subsection{Minimum-image torus sequence}

For the simulation convention, take even $L$ and let $D_L$ be the centered
square of minimum-image representatives.  Define
\begin{equation}
 J_L(R)=\frac{\cL}{|R|^4},\qquad
 \cL=\frac4{\displaystyle
 \sum_{R\in D_L\setminus\{0\}}|R|^{-4}}.               \label{eq:MI}
\end{equation}
For the infinite lattice,
\begin{equation}
 \cinf=
 \frac4{\displaystyle
 \sum_{R\in\mathbb Z^2\setminus\{0\}}|R|^{-4}}
 =
 \frac6{\pi^2G}
 =
 0.6637008046\ldots,                                   \label{eq:c-infinity}
\end{equation}
where $G$ is Catalan's constant.  The omitted tail is $O(L^{-2})$, and
normalization of both interactions to total weight four gives
\begin{equation}
 \|J_L^0-J_\infty\|_{\ell^1}=O(L^{-2}),                \label{eq:l1}
\end{equation}
where $J_L^0$ is extended by zero outside $D_L$.  Therefore
\begin{equation}
 \sup_q|E_L(q)-E_\infty(q)|
 \le2\|J_L^0-J_\infty\|_{\ell^1}
 =O(L^{-2}).                                           \label{eq:uniform-MI}
\end{equation}
Assumption~\ref{ass:uniform} and the marginal infrared divergence both hold.

\begin{corollary}[Two-dimensional minimum-image marginal model]
\label{cor:MI}
For the sequence \eqref{eq:MI}, every finite $n\ge2$, and every $T>0$,
\begin{equation}
 \lim_{L\to\infty}
 \langle|\bm M_L|^2\rangle_{L,0}=0.                   \label{eq:MI-result}
\end{equation}
\end{corollary}

\subsection{Quantitative finite-size envelope at marginality}

The same finite-volume ingredients give more than the limiting statement in
Corollary~\ref{cor:MI}.  They quantify the doubly logarithmic obstruction
directly for the zero-field observable.

\begin{lemma}[Marginal reciprocal sum]
\label{lem:quantitative-sum}
For $h_L=T/L^2$, define
\begin{equation}
 A_L(h_L)=\frac1{L^2}\sum_{q\in\mathcal Q_L}
 \frac1{h_L+E_L(q)}.
\label{eq:quantitative-A}
\end{equation}
For every fixed $T>0$, there are constants $a_T>0$ and $L_T<\infty$ such
that, for every even $L\ge L_T$,
\begin{equation}
 A_L(T/L^2)\ge a_T\log\log L.
\label{eq:quantitative-sum}
\end{equation}
\end{lemma}

\begin{proof}
Equations~\eqref{eq:marginal-upper} and \eqref{eq:uniform-MI} imply that
there are constants $B,C,K>0$ such that, for all sufficiently large even
$L$ and $0<|q|\le1/2$,
\begin{equation}
 E_L(q)\le \frac{B}{L^2}+C|q|^2\log\frac K{|q|}.
\label{eq:finite-L-kernel-upper}
\end{equation}
Represent momenta by $q_k=2\pi k/L$ with $k\in D_L$.  Fix
$0<\alpha<(4\pi\sqrt2)^{-1}$ and let
\begin{equation}
 \mathcal S_m=\{k\in D_L:\|k\|_\infty=m\},
 \qquad1\le m\le\lfloor\alpha L\rfloor.
\end{equation}
Every such shell contains exactly $8m$ points, and its momenta satisfy
$2\pi m/L\le|q_k|\le2\pi\sqrt2m/L\le1/2$.  After enlarging a constant
$D_T<\infty$, Eq.~\eqref{eq:finite-L-kernel-upper} gives uniformly on these
shells
\begin{equation}
 h_L+E_L(q_k)
 \le\frac{D_T}{L^2},
 m^2\log\frac{K'L}{m}.
\end{equation}
Discarding the nonnegative zero mode and summing the remaining shells yields
\begin{align}
 A_L(h_L)
 &\ge\frac8{D_T}
 \sum_{m=1}^{\lfloor\alpha L\rfloor}
 \frac1{m\log(K'L/m)} \nonumber\\
 &\ge\frac8{D_T}
 \left[
  \log\log(K'L)
  -\log\log\!\left(\frac{2K'}{\alpha}\right)
 \right].
\label{eq:shell-loglog}
\end{align}
The last line follows by comparison with the integral of
$[x\log(K'L/x)]^{-1}$; $K'$ may be enlarged so that this function is
decreasing throughout the summation range.  Equation~\eqref{eq:quantitative-sum}
follows for all sufficiently large $L$.  The omitted zero mode is exactly
$L^{-2}h_L^{-1}=1/T$, so it contributes no logarithmic growth.
\end{proof}

\begin{corollary}[Quantitative absence of vector LRO]
\label{cor:quantitative-no-lro}
For every finite $n\ge2$ and every fixed $T>0$, there are
$C_{n,T},L_{n,T}<\infty$ such that, for every even $L\ge L_{n,T}$,
\begin{equation}
 \boxed{\displaystyle
 \langle|\bm M_L|^2\rangle_{L,0}
 \le\frac{C_{n,T}}{\sqrt{\log\log L}}.}
\label{eq:quantitative-M2}
\end{equation}
\end{corollary}

\begin{proof}
Apply Eq.~\eqref{eq:master-bound} at $h_L=T/L^2$ and use
Lemma~\ref{lem:quantitative-sum}:
\begin{equation}
 m_{L,h_L}
 \le\frac1{\sqrt{(n-1)T A_L(h_L)}}.
\end{equation}
The finite-volume bridge \eqref{eq:bridge} is valid for this $L$-dependent
field.  Since $\beta h_LL^2=1$, it gives
\begin{equation}
 \langle|\bm M_L|^2\rangle_{L,0}
 \le\frac{n}{\tanh(1)}m_{L,h_L}.
\end{equation}
Combining the last two displays with Eq.~\eqref{eq:quantitative-sum} proves
Eq.~\eqref{eq:quantitative-M2}.
\end{proof}

Equation~\eqref{eq:quantitative-M2} is a rigorous asymptotic upper envelope.
Its constants are not optimized, and it is neither a matching lower bound nor
an exact finite-size decay law.

\subsection{The boundary of the criterion}

For Eq.~\eqref{eq:power-law} with $0<\sigma<2$,
\begin{equation}
 E(q)\asymp |q|^\sigma,\qquad
 \int_0^\delta\frac{k\,\dd k}{k^\sigma}<\infty.        \label{eq:sigma-lt2}
\end{equation}
Theorem~\ref{thm:main} is silent in this regime, and continuous-symmetry LRO
is not excluded.  For $\sigma=2$, the integral grows as $\log\log L$; for
$\sigma>2$, the leading analytic kernel is $q^2$ and the growth is
$\log L$.  Figure~\ref{fig:ir-regimes} displays these three cases using
their continuum proxy kernels.

\begin{figure*}[t]
 \centering
 \includegraphics[width=0.94\textwidth]
 {paper/figures/on_infrared_regimes.pdf}
 \caption{\label{fig:ir-regimes}
 Two-dimensional infrared regimes represented by deterministic continuum
 proxies.  (a) For $\sigma<2$, $E_\sigma(k)=k^\sigma$; at the marginal
 point $E_2(k)=k^2\log(e/k)$; above it the leading analytic form is $k^2$.
 (b) The radial reciprocal integral with $k_{\min}=2\pi/L$ approaches a
 finite limit for $\sigma=1.75$, grows as $\log\log L$ for $\sigma=2$, and
 as $\log L$ for $\sigma=2.25$.  These curves illustrate the theorem's
 infrared hypothesis; they are neither Monte Carlo measurements nor proof
 premises.}
\end{figure*}

\section{What changes from XY to \texorpdfstring{$O(n)$}{O(n)}}
\label{sec:on-vs-xy}

Table~\ref{tab:on-comparison} summarizes the points at which the
generalization is substantive.  The sphere geometry introduces no extra
term because the chosen derivatives preserve the single-site measure.
There is also no odd--even obstruction in the present proof.  Each
$(1,a)$ plane is treated separately, and the inequalities are summed only
after they are established.  Odd--even effects in some soft-spin
fluctuation spectra concern a different approximation and do not alter
Eq.~\eqref{eq:master-bound} \cite{Nussinov2001}.

\begin{table*}[t]
\caption{\label{tab:on-comparison}
The XY proof and its hard-spin $O(n)$ extension.  The theorem has a common
no-LRO conclusion, but it does not assert that the remaining phases have
the same correlations or topological defects.}
\centering
\renewcommand{\arraystretch}{1.18}
\resizebox{0.98\textwidth}{!}{%
\begin{tabular}{lll}
\toprule
Step & XY, $n=2$ & General hard-spin $O(n)$\\
\midrule
Single-site calculus &
\shortstack[l]{one periodic angle and\\$\partial_{\theta_x}$} &
\shortstack[l]{sphere rotations generated by\\$\mathcal L_x^{1a}$}\\
\addlinespace[4pt]
Measure identity &
periodic integration by parts &
\shortstack[l]{invariance of surface measure\\under $SO(n)$}\\
\addlinespace[4pt]
Pair second variation &
the planar spin scalar product &
\shortstack[l]{$S_x^1S_y^1+S_x^aS_y^a$,\\a selected-plane projection}\\
\addlinespace[4pt]
Transverse budget &
\shortstack[l]{one channel,\\$V\sum_x\sin^2\theta_x$} &
\shortstack[l]{$n-1$ channels with total budget\\
$V\sum_x[1-(S_x^1)^2]$}\\
\addlinespace[4pt]
Zero-field tensor &
\shortstack[l]{$\langle(M^1)^2\rangle=$\\$\langle|\bm M|^2\rangle/2$} &
\shortstack[l]{$\langle(M^1)^2\rangle=$\\$\langle|\bm M|^2\rangle/n$}\\
\addlinespace[4pt]
Possible no-order physics &
\shortstack[l]{BKT and marginal logarithmic\\scenarios are relevant} &
\shortstack[l]{no common BKT or correlation-law\\conclusion follows for $n\ge3$}\\
\bottomrule
\end{tabular}
}
\end{table*}

The hard-spin constraint is used twice.  It bounds the internal-plane
projection in Eq.~\eqref{eq:projection-bound}, and it makes
$M_L^1\in[-1,1]$ in Lemma~\ref{lem:tilt}.  A soft-spin $\phi^4$ theory may
obey a related no-order theorem, but it requires moment control and a
modified tilt argument.  Likewise, a general compact homogeneous space
would require enough divergence-free symmetry generators and a
model-specific second-variation bound.  Those extensions cannot be inferred
from the notation $O(n)$ alone.

The spatial kernel itself does not depend on $n$.  Once $J_L$ is fixed,
the deterministic calculation of $E_L(q)$ is common to XY, Heisenberg, and
higher-component spins.  The component count enters the proof through the
factor $n-1$ and through the isotropic zero-field factor $1/n$.  No
component-by-component Monte Carlo campaign is needed to verify the
Fourier-kernel premise.

\section{Deterministic lattice audit}
\label{sec:kernel-audit}

For the normalized square lattice,
\begin{equation}
 \rhoLog=\frac{\pi\cinf}{2}
 =1.042538785978\ldots.                                \label{eq:rho-log}
\end{equation}
We evaluated the exact minimum-image sum with
\begin{equation}
 1-\cos(q\cdot R)=2\sin^2\!\left(\frac{q\cdot R}{2}\right)             \label{eq:stable-phase}
\end{equation}
to avoid small-angle cancellation.  Sizes range from $L=32$ to $8192$,
and the four fixed momentum-index paths are
$(1,0)$, $(2,0)$, $(1,1)$, and $(2,1)$.

Let
\begin{equation}
 Y_L(\bm n)=
 \frac{E_L(2\pi\bm n/L)}{|2\pi\bm n/L|^2},
\end{equation}
and define the intercept-free dyadic slope
\begin{equation}
 \rho_{\mathrm{eff}}(L;\bm n)
 =
 \frac{Y_{2L}(\bm n)-Y_L(\bm n)}{\log2}.                \label{eq:rho-effective}
\end{equation}
At the largest dyadic pair, the four relative discrepancies from
\eqref{eq:rho-log} lie between $8.16\times10^{-7}$ and
$9.05\times10^{-7}$.  Direct use of $1-\cos$ and the stable
formula \eqref{eq:stable-phase} differs by $2.38\times10^{-12}$ relative at
$L=8192$.

\begin{figure*}[t]
 \centering
 \includegraphics[width=0.96\textwidth]
 {paper/figures/kernel_audit.pdf}
 \caption{\label{fig:kernel}
 Deterministic audit of the exact minimum-image sequence.
 (a) $E_L(\bm q)/|\bm q|^2$ has the same logarithmic slope on four momentum
 paths.  (b) The intercept-free slope
 \eqref{eq:rho-effective} converges to $\rhoLog$ with a final relative
 discrepancy below $9.1\times10^{-7}$.
 (c) The $L^2$-scaled normalization corrections remain bounded; the
 certified uniform-kernel comparison has the expected relative
 $O(1/\log L)$ scale at $\kmin$.  The calculation audits the model
 convention and implementation.  Corollary~\ref{cor:MI} is analytic.}
\end{figure*}

The minimum-image omission changes $E_L(\kmin)$ by $O(L^{-2})$, while the
marginal kernel itself is $O(L^{-2}\log L)$.  Thus the relative convention
effect is only $O(1/\log L)$.  It vanishes in the thermodynamic limit but is
of the same form used in finite-size intercept extrapolations.

\section{The \texorpdfstring{$n=2$}{n=2} finite-size application}
\label{sec:xy-application}

The theorem settles only the existence of a positive thermodynamic
magnetization.  It does not fix the asymptotic correlation law at
$\sigma=2$.  To see why accessible XY data can appear ordered, consider the
scaling family
\begin{equation}
 C_L(Lx)=A[\ell+\phi(x)]^{-p},\qquad
 \ell=\log(L/L_0),\quad p>0.                           \label{eq:log-family}
\end{equation}
Expanding at large $\ell$,
\begin{equation}
 C_L(Lx)=A\ell^{-p}
 \left[1-\frac{p\phi(x)}{\ell}+O(\ell^{-2})\right].    \label{eq:log-expansion}
\end{equation}
The constant spatial term survives at zero momentum and cancels at a
nonzero torus momentum.  Consequently,
\begin{align}
 \langle|\bm M_L|^2\rangle
 &\sim A\ell^{-p},                                      \label{eq:M2-log}\\
 \langle|\bm M_{\kmin}|^2\rangle
 &\sim Bp\,\ell^{-p-1},                                \label{eq:Mk-log}\\
 \chi_{\kmin}
 &\sim BpL^2\ell^{-p-1},                               \label{eq:chi-log}\\
 \left(\frac{\xi}{L}\right)^2
 &\sim\frac{\ell}{p}.                                  \label{eq:xi-log}
\end{align}
For small $p$, Eqs.~\eqref{eq:chi-log}--\eqref{eq:xi-log} look like
$L^2/\log L$ and $\log L$, while Eq.~\eqref{eq:M2-log} changes by only a few
percent.

The public Zenodo record \cite{YaoData2025} provides run-level means and
reported uncertainties, not synchronized bin means or empirical
cross-observable covariance.  We therefore treat the following calculations
as identifiability diagnostics, not as a replacement for a covariance-aware
reanalysis of the underlying measurements.  On 32 registered
$(\beta,L_{\min})$ windows, AICc selects the logarithmic model in 29, a BKT
model in two, and an ordered model in one.  Profiling a common correlation
coefficient over $[-0.95,0.95]$ leaves 29 windows logarithmic-only, two
logarithmic/ordered ambiguous, and one ordered-only.

\begin{figure*}[t]
 \centering
 \includegraphics[width=0.93\textwidth]
 {paper/figures/window_comparison.pdf}
 \caption{\label{fig:windows}
 Registered finite-size model comparison for the public XY summaries.
 (a) AICc winners across inverse temperatures and lower fit cutoffs.
 (b) Signed AICc separation between the BKT and logarithmic candidates.
 The changing window contains genuine model-selection information: a
 positive intercept obtained from a single cutoff is not an asymptotic
 lower bound.  Missing synchronized covariance prevents this figure from
 serving as a definitive statistical test.}
\end{figure*}

An exact Gaussian compact-phase calculation supplies a control with a known
zero asymptote and the same marginal kernel.  Fitting those sequences to an
ordered intercept form gives positive intercepts in all 12 combinations of
four temperatures and three lower cutoffs.  The fitted values range from
$0.504$ to $0.923$, despite the analytically known zero limit.

\begin{figure*}[t]
 \centering
 \includegraphics[width=0.93\textwidth]
 {paper/figures/gaussian_false_intercept.pdf}
 \caption{\label{fig:gaussian}
 Known-zero Gaussian control using the marginal kernel.
 (a) The squared magnetization varies slowly over the simulated size range
 while tending to zero analytically.  (b) An ordered-intercept fit returns a
 positive $\widehat g_0$ in all 12 registered temperature/window
 combinations.  The control demonstrates finite-window false positives; it
 is not a claim that vortices or nonlinear effects are Gaussian.}
\end{figure*}

Synthetic recovery tests quantify the remaining overlap.  For
$L_{\min}=256$, $L_{\max}=4096$, and the reported error scale, the correct
recovery probabilities for logarithmic and BKT truths are only
$0.738$--$0.829$ at $\beta=4,8$, whereas ordered truths are recovered with
probability $0.935$--$0.970$.  This asymmetry explains why an
ordered-looking plateau is easier to identify than an exceptionally slow
zero asymptote.

\begin{figure*}[t]
 \centering
 \includegraphics[width=0.93\textwidth]
 {paper/figures/identifiability_covariance.pdf}
 \caption{\label{fig:identifiability}
 Sensitivity to unavailable covariance and finite-window identifiability.
 (a) Winner sets obtained by profiling the registered covariance range.
 (b) Synthetic recovery at $L_{\min}=256$ and $L_{\max}=4096$; the dashed
 line marks a 90\% target.  Slow no-order alternatives remain below that
 target at the public precision.  Synchronized bins or jackknife replicas
 are required for a final joint likelihood.}
\end{figure*}

The theorem and these diagnostics answer different questions.
Theorem~\ref{thm:main} excludes a positive thermodynamic intercept under its
stated model hypotheses.  The fits explain how a finite data set can remain
consistent with such an intercept.  They do not establish whether the
ultimate $n=2$ correlations are logarithmic or algebraic.

\section{Relation to rigorous results}
\label{sec:literature}

Mermin's classical inequality \cite{Mermin1967} provides the conceptual
ancestor of Eqs.~\eqref{eq:gibbs-ibp-one}--\eqref{eq:classical-bog}.
Nussinov \cite{Nussinov2001} formulates classical $O(n)$ rotations with
\begin{equation}
 L_{\alpha\gamma}
 =
 S_\alpha\partial_{S_\gamma}
 -
 S_\gamma\partial_{S_\alpha}
\end{equation}
and derives the invariant-measure integration by parts.  Some displayed
steps there are specialized to $n=2$; the projected pair variation and
the combined $n-1$ transverse budget in
Eqs.~\eqref{eq:pair-variation}--\eqref{eq:parseval-budget} make the present
finite-volume statement explicit for general hard spins.

Bruno's theorem \cite{Bruno2001} is a quantum-spin magnetization result.
Its Fourier criterion includes the marginal $d=2$, $\alpha=4$ point, but no
formal $S\to\infty$ limit is needed here.  Messager,
Miracle-Sol\'e, and Ruiz \cite{Messager1984} obtain pointwise correlation
bounds for fixed $SO(n)$-symmetric long-range systems.  Ioffe, Shlosman,
and Velenik \cite{Ioffe2002} prove invariance of Gibbs states for recurrent
interaction walks under broad compact-symmetry hypotheses.  For
$J(R)\propto r^{-4}$,
\begin{equation}
 1-\widehat p(q)\propto q^2\log(1/q),
\end{equation}
and the associated walk is recurrent.  These results support the same
no-symmetry-breaking conclusion in their respective infinite-volume
settings.

The contribution here is narrower in microscopic scope and more direct in
finite-volume observables.  It permits an explicitly size-dependent
minimum-image interaction sequence, keeps the physical order of limits, and
ends at the zero-field squared magnetization without invoking an
infinite-volume ergodic decomposition.

\section{Limitations and open questions}
\label{sec:limitations}

The theorem proves absence of vector LRO.  It does not prove exponential
clustering, a mass gap, a universal logarithmic exponent, or a common phase
diagram for all $n$.  For $n=2$, a harmonic treatment of
Eq.~\eqref{eq:intro-marginal} suggests
\begin{equation}
 C(r)\sim[\log r]^{-p(T)},                              \label{eq:spin-wave-log}
\end{equation}
but the compact field also contains vortices.  Renormalization-group
analyses around a short-range BKT line give a local long-range eigenvalue
\begin{equation}
 \Delta_g=2-\sigma-\eta_{\rm sr},                       \label{eq:rg-eigenvalue}
\end{equation}
which is negative at $\sigma=2$ when $\eta_{\rm sr}>0$
\cite{Giachetti2021}.  Recent coupled spin-wave--vortex work likewise
supports persistence of BKT behavior \cite{Walther2026}.  Neither statement
is an exact global theorem selecting the eventual correlation law.

For $n\ge3$, the XY vortex mechanism does not carry over unchanged.  The
finite-volume inequality rules out magnetization but supplies no
$n$-independent replacement for the BKT description.  Establishing
clustering or a correlation-length law requires additional methods.

A soft-spin extension is plausible but not automatic.  One would need
uniform control of the unbounded radial moments, a denominator estimate that
survives the radial potential, and a replacement for the bounded tilt
Lemma~\ref{lem:tilt}.  A target-manifold extension would require a family of
measure-preserving Killing fields whose squared actions control the chosen
order parameter.  Frustrated interactions would require an absolute-coupling
kernel or a different denominator bound.  These are distinct theorems, not
corollaries of Theorem~\ref{thm:main}.

\section{Conclusion}

The classical hard-spin $O(n)$ generalization closes without an XY-specific
failure.  Its nontrivial ingredients are the invariant rotation calculus on
$S^{n-1}$, the selected-plane form of the pair second variation, and the
shared Parseval budget of all $n-1$ transverse modes.  Together they give the
finite-volume estimate \eqref{eq:master-bound}.  Uniform kernel convergence
and the regulated infrared divergence eliminate a field-selected
magnetization, while the bounded exponential tilt transfers that conclusion
to the zero-field squared magnetization.

In two dimensions, finite-second-moment interactions and the marginal
$1/r^4$ interaction satisfy the infrared hypothesis.  The precise
minimum-image, size-normalized torus sequence converges uniformly to that
fixed model, so its thermodynamic squared magnetization vanishes for every
finite $n\ge2$ and $T>0$.  More quantitatively, the zero-field squared
magnetization is bounded above by
$C_{n,T}/\sqrt{\log\log L}$ for all sufficiently large even $L$.  This
envelope makes the marginal obstruction explicit without claiming a sharp
rate at accessible sizes.  The deterministic kernel calculation confirms
the model convention and exact logarithmic coefficient; the public-data
tests show why the XY finite-size sequence can nevertheless resemble a
positive plateau.  The result is a generic theorem for the stated class of
classical $O(n)$ spin systems, not a theorem for every two-dimensional
continuous field theory.

\begin{acknowledgments}
Acknowledgments, funding information, and the final author list will be
added before submission.
\end{acknowledgments}

\appendix

\section{Matrix form of the pair identity}
\label{app:matrix}

Let $G^{1a}$ be the antisymmetric $n\times n$ matrix with
\begin{equation}
 G^{1a}_{1a}=-1,\qquad G^{1a}_{a1}=1,
\end{equation}
and all other entries zero.  For a column vector $s$,
$G^{1a}s$ is tangent to $S^{n-1}$ and
\begin{equation}
 (G^{1a})^2=-P_{1a},
\end{equation}
where $P_{1a}$ projects onto the $(1,a)$ plane.  For
$f(s,t)=-J\,s^\mathsf Tt$,
\begin{align}
 \partial_s^2f
 &=-J[(G^{1a})^2s]^\mathsf Tt
 =J\,s^\mathsf TP_{1a}t,\\
 \partial_t^2f
 &=-Js^\mathsf T(G^{1a})^2t
 =J\,s^\mathsf TP_{1a}t,\\
 \partial_s\partial_tf
 &=-J(G^{1a}s)^\mathsf T(G^{1a}t)
 =-J\,s^\mathsf TP_{1a}t.
\end{align}
The Hermitian Fourier combination is therefore
\begin{align}
 &|u|^2\partial_s^2f+\bar uv\,\partial_s\partial_tf
 +\bar vu\,\partial_t\partial_sf+|v|^2\partial_t^2f\nonumber\\
 &\qquad =
 J|u-v|^2s^\mathsf TP_{1a}t,
\end{align}
which is Eq.~\eqref{eq:pair-variation}.  This derivation also shows that
components orthogonal to the selected plane do not enter the second
variation.

\section{A covariance proof used in the tilt lemma}
\label{app:covariance}

For a real random variable $Q$ and nondecreasing functions $f$ and $g$, let
$Q'$ be an independent copy.  Then
\begin{align}
 2\operatorname{Cov}(f(Q),g(Q))
 &=
 \Exp\!\left\{[f(Q)-f(Q')]\right.\nonumber\\
 &\qquad\left.\times[g(Q)-g(Q')]\right\}\ge0.           \label{eq:chebyshev}
\end{align}
Taking $Q=|X|$, $f(Q)=Q^2$, and $g(Q)=\cosh(tQ)$ proves the covariance step
in Lemma~\ref{lem:tilt}.  Compactness enters through $Q\le1$ in the
concavity bound $\tanh(tQ)\ge Q\tanh(t)$.

\section{Reproducibility record}
\label{app:reproducibility}

The proof identities are checked by deterministic scripts for
$n=2,3,4,8$.  The maximum residual in the pair second-variation identity is
$2.7\times10^{-15}$, and the discrete transverse Parseval identity is
verified to floating-point precision.  Negative projected-pair samples are
included to guard against replacing the averaged positivity
\eqref{eq:gibbs-ibp-two} by an incorrect pointwise statement.

The committed machine-readable records are stored as
\begin{center}
\begin{minipage}{0.92\columnwidth}
\ttfamily
artifacts/on\_theorem\_audit.json\\
artifacts/on\_infrared\_regimes.json\\
artifacts/kernel\_audit.json
\end{minipage}
\end{center}
The figure generator reads only these versioned records and writes vector
PDFs with fixed metadata.  The numerical certificates are regression checks
of algebra, conventions, and implementation; none is substituted for an
analytic step in Theorem~\ref{thm:main}.

\bibliography{paper/references}

\end{document}
